\begin{problem}[第32届全国中学生物理竞赛复赛][多普勒频移]{43}
    如图，飞机在距水平地面（$xz$ 平面）等高的航线 $KA$（沿 $x$ 正方向）上，以大小为 $v$（$v$ 远小于真空中的光速 $c$）的速度匀速飞行；机载雷达天线持续向航线正右侧地面上的被测固定目标 $P$ 点（其 $x$ 坐标为 $x_{\mathrm{P}}$）发射扇形无线电波束（扇形的角平分线与航线垂直），波束平面与水平地面交于线段 $BC$（$BC$ 随着飞机移动，且在测量时应覆盖被测目标 $P$ 点），取 $K$ 点在地面的正投影 $O$ 为坐标原点。已知 $BC$ 与航线 $KA$ 的距离为 $R_{0}$。天线发出的无线电波束是周期性的等幅高频脉冲余弦波，其频率为 $f_{0}$。
    \begin{subproblem}
        已知机载雷达天线经过 $A$ 点（其 $x$ 坐标为 $x_{\mathrm{A}}$）及此后朝 $P$ 点相继发出无线电波信号，由 $P$ 反射后又被机载雷达天线接收到，求接收到的回波信号的频率与发出信号的频率之差（频移）。
    \end{subproblem}
    \begin{subproblem}
        已知 $BC$ 长度为 $L_{s}$，讨论上述频移分别为正、零或负的条件，并求出最大的正、负频移。
    \end{subproblem}
    \begin{subproblem}
        已知 $R_{0} >> L_{s}$，求从 $C$ 先到达 $P$ 点、直至 $B$ 到达 $P$ 点过程中最大频移与最小频移之差（带宽），并将其表示成扇形波束的张角 $\theta$ 的函数。
    \end{subproblem}
\end{problem}